\section*{Chapter \ref{ch:g6:describing-motion} --- Describing Motion: Trajectory and Speed}

\begin{solution}{exo:g6:describing-motion:1}
Straight (down the pane); circular; curved (an arch); curved (a
weaving S-line).
\end{solution}

\begin{solution}{exo:g6:describing-motion:2}
The hub: a straight line, gliding parallel to the road. The valve:
a curve of graceful hops --- rising, arching, touching down with
each turn of the wheel.
\end{solution}

\begin{solution}{exo:g6:describing-motion:3}
Uniform; varied (speeding up); uniform; varied (slowing down).
\end{solution}

\begin{solution}{exo:g6:describing-motion:4}
Equal gaps: \omterm{def:g6:describing-motion:uniform}{uniform motion}. Growing gaps ($1, 2, 4, 7$): speeding
up --- more distance in each equal tick.
\end{solution}

\begin{solution}{exo:g6:describing-motion:5}
A drifting camera adds its own motion to the record --- the dots
then mix the object's journey with the camera's. Motion is always
motion compared to something, and the film must fix its
``something'': the remark on viewpoints.
\end{solution}

\begin{solution}{exo:g6:describing-motion:6}
Table: $30$ min $\to$ \qty{12}{km}, so $60$ min $\to$
\qty{24}{km}: the ferry's \omterm{def:g5:speed-distance-time:speed}{speed} is $24$ kilometres per hour.
\end{solution}

\begin{solution}{exo:g6:describing-motion:7}
$60$ min $\to$ \qty{14}{km}; $30$ min $\to$ \qty{7}{km}; $15$ min
$\to$ \qty{3.5}{km}.
\end{solution}

\begin{solution}{exo:g6:describing-motion:8}
A stroller at $12$ kilometres per hour: rejected --- three times
the strolling pace. The bus: $10$ km in $30$ min is $20$
kilometres per hour --- exactly city-bus pace, accepted. The
train: $75$ km in $15$ min is $300$ kilometres per hour ---
top-of-the-board but plausible for a \omterm{def:g5:speed-distance-time:speed}{high-speed} train: accepted.
\end{solution}

\begin{solution}{exo:g6:describing-motion:9}
For the passenger opposite: no motion at all --- no \omterm{def:g6:describing-motion:trajectory}{trajectory},
pace zero. For the cow: a long \omterm{def:g6:describing-motion:trajectory}{straight trajectory} at the train's
full pace. Both honest: each describes the dozer compared to their
own viewpoint, and motion is always motion compared to something.
\end{solution}

\begin{solution}{exo:g6:describing-motion:10}
Riding time: $18 - 6 = 12$ minutes for \qty{9}{km}. Table: $12$
min $\to$ $9$ km, so $60$ min $\to$ $45$ km --- exactly $45$
kilometres per hour: neither above nor below, but spot on.
\end{solution}

\begin{solution}{exo:g6:describing-motion:11}
Film the fetch with a still camera; pause every half-second; mark
the dog's positions against fence posts; \omterm{def:g6:measurement-in-science:measuring}{measure} the gaps. Even
gaps throughout: steady pace. Gaps that spread, then shrink (or
alternate wide and narrow): sprint-and-coast. The dots deliver the
verdict, whatever the family shouted.
\end{solution}

\begin{solution}{exo:g6:describing-motion:12}
From the room: a curve --- the ball arches forward and down from
the table edge to the floor. From the little cart rolling at the
ball's forward pace: the ball keeps no lead and no lag --- it
appears to fall \emph{straight down}. Same fall, two viewpoints,
and the simpler view belongs to the moving observer: a seed of a
great idea.
\end{solution}

\begin{solution}{pb:g6:describing-motion:1}
\textbf{1.} Equal gaps of \qty{50}{m}: \omterm{def:g6:describing-motion:uniform}{uniform motion}.
\textbf{2.} $5 \times 50 = \qty{250}{m}$.
\textbf{3.} The tram pulls away from rest: short gaps growing ---
speeding up over the first three ticks --- then settling to
steady \qty{50}{m} gaps: cruising.
\textbf{4.} The last two intervals, at \qty{50}{m} each.
\textbf{5.} $10$ s $\to$ \qty{50}{m}, so $60$ s $\to$
\qty{300}{m} per minute.
\textbf{6.} $60$ min $\to$ $60 \times 300 = \qty{18000}{m}$:
cruising \omterm{def:g5:speed-distance-time:speed}{speed} $18$ kilometres per hour.
\textbf{7.} Table: \qty{18}{km} per $60$ min, so \qty{12}{km} in
$40$ minutes.
\textbf{8.} $50 - 40 = 10$ minutes lost to stops and slow zones.
\textbf{9.} $50$ min $\to$ \qty{12}{km}, so $60$ min $\to$
$12 \times 60 \div 50 = \qty{14.4}{km}$: the whole journey runs at
about $14$ kilometres per hour --- the poster's $18$ is honest
only for the cruising stretches, not ``clear across town''.
\textbf{10.} A whole-journey \omterm{def:g5:speed-distance-time:speed}{speed} describes the trip \emph{as if}
it were uniform: the one steady pace that would cover the same
\qty{12}{km} in the same $50$ minutes. No speedometer shows it,
yet it is the fair number for planning, because it already
contains every stop and slow zone.
\textbf{11.} $50 - 4 = 46$ minutes; $12 \times 60 \div 46 \approx
\qty{15.7}{km}$ per hour --- call it about $16$.
\textbf{12.} For example: ``Commuters should plan by the
whole-journey \omterm{def:g5:speed-distance-time:speed}{speed} --- about $14$ kilometres per hour today ---
because it is the only number that includes the waiting; the
cruising $18$ flatters the line and misses every stop.''
\end{solution}
